\section*{Probability Spaces}

A probability space is a mathematical model used to represent a random experiment. It consists of a sample space, a $\sigma$-algebra, and a probability measure. The sample space is the set of all possible outcomes of the experiment. During the random experiment, one outcome is selected, which represents a realization of the experiment. We may not be able to observe the chosen outcome directly, but only whether it belongs to a subset of the sample space. We may also be interested in whether two events occur simultaneously or whether at least one of them happens. The theory of probability must be rich enough to combine events using basic set operations.

To model the random experiment, we assign probabilities to the events in a consistent manner, which is captured by the probability measure function. The sample space can be a metric space or an abstract space. In this chapter, we first introduce measure and probability spaces in general and go through some basic properties. Then, we focus on sample spaces in which we can talk about open sets, such as Euclidean space and metric space, and discuss the Borel algebra generated by the open sets. This will allow us to study probability measures on more specific sample spaces.

\subsection*{2.1 Countable Sets}

A set is countable if the elements can be listed sequentially.

\begin{defbox}{2.1}
Two sets are said to have the same size, or cardinality, if there exists a bijection between them. A set is said to be \textit{countable} if it can be mapped bijectively to the set of natural numbers. An infinite set that is not countable is called \textit{uncountable}.
\end{defbox}

\textbf{Example 2.1.1 (Countable Sets)} \\
\begin{itemize}
    \item The set of positive even integers is a subset of natural numbers $\mathbb{N}$. Nevertheless, it has the same cardinality as $\mathbb{N}$ and is therefore countable. We have an explicit bijection $f(x)=2x$ from $\mathbb{N}$ onto the positive even integers.
    \item The set of integers $\mathbb{Z}$ are countable because we can list the integers as follows: $0,1,-1,2,-2,3,-3,\ldots$
    \item An infinite subset of a countable set is countable.
    \item The union of two countable sets is countable. In particular, if an uncountable set is the union of two sets, then one of these sets must be uncountable.
    \item The Cartesian product of two countable sets is countable.
    \item The set of positive rational numbers is countable because it can be put into one-to-one correspondence with the subset
    \[
    \{(a,b)\in \mathbb{N}\times \mathbb{N} : \gcd(a,b)=1\},
    \]
    which is a subset of $\mathbb{N}\times \mathbb{N}$. The set of rational numbers between $0$ and $1$ is also countable because it is a subset of the positive rational numbers.
\end{itemize}

\noindent$\blacktriangleright$ \textbf{Convention} Some people extend the notion of countable set to include finite sets as well. In this book, however, we will use the term ``countable'' exclusively for infinite sets that are in bijection with the set of natural numbers. To refer to sets that are either finite or countable, we will use the term \textit{at most countable}. We say that a set is \textit{infinitely countable} if we want to emphasize that the set has infinite cardinality.

\textbf{Example 2.1.2 (Cantor's Diagonal Argument)} \\
Using Cantor's diagonal argument, we can show that the set of real numbers in the unit interval $[0,1]$ is uncountable. To do this, we first consider the set $S$ of infinite binary sequences. Each such sequence has the form $(b_1,b_2,b_3,b_4,\dots)$, with $b_k\in \{0,1\}$ for all $k\ge 1$.

To prove that the set $S$ is uncountable, we assume for the sake of contradiction that $S$ is countable and hence can be enumerated as a list containing all the infinite binary sequences. Then we can enumerate the binary sequences in a list. This can be viewed as an infinite array of binary numbers. We denote the $j$-th bit in the $i$-th sequence by $b_{ij}$. We assume that the list is complete.

We then construct another binary sequence $s^{*}$ by extracting and flipping the diagonal entries in the array of bits, starting with the first bit of the first sequence, the second bit of the second sequence, and so on. The $i$-th bit in $s^{*}$ is defined as $1-b_{ii}$. Since $s^{*}$ differs with each infinite sequence in the list in at least one position, it is not equal to any one of them. This contradicts the assumption that we can exhaustively enumerate all the infinite binary sequences as a list.

For example, suppose we have a list of binary sequences as follows:
\[
s_1 = 1011011000111001\ldots
\]
\[
s_2 = 1111000001101101\ldots
\]
\[
s_3 = 0010000001001101\ldots
\]
\[
s_4 = 1010001110000111\ldots
\]
\[
\vdots
\]
Each sequence in the list has infinitely many digits, and the $i$-th bit in the $i$-th number in the list is displayed in boldface. In this case, we have $s^{*}=0001\ldots$, which is not in the list.

Next, we associate a binary sequence $(b_1,b_2,b_3,\ldots)$ with a real number by the formula
\[
\sum_{j=1}^{\infty} b_j 2^{-j}.
\]

However, this map is not necessarily injective, as different sequences can correspond to the same real number. For example, $0.10000\ldots$ and $0.011111\ldots$ are both equal to $1/2$. To obtain an injection, we consider only the binary sequences that do not end with infinitely many trailing zeros. The removed sequences correspond to the dyadic rationals in $[0,1]$, which are countable. Since we have removed only a countable set from $S$, the resulting set of binary sequences that do not end with infinitely trailing zeros is uncountable. This means that there exists a one-to-one function from an uncountable set to the interval $[0,1]$. Hence, $[0,1]$ is uncountable.

One can prove that the set of all real numbers and the interval $(0,1)$ has the same cardinality. For example, we can use the arctangent function $\frac{1}{\pi}\tan^{-1}(x)+0.5$. Therefore, the set of real numbers is uncountable.
